% Options for packages loaded elsewhere
\PassOptionsToPackage{unicode}{hyperref}
\PassOptionsToPackage{hyphens}{url}
\PassOptionsToPackage{dvipsnames,svgnames,x11names}{xcolor}
\documentclass[
  10pt,
]{article}
\usepackage{xcolor}
\usepackage[margin=2.4cm]{geometry}
\usepackage{amsmath,amssymb}
\setcounter{secnumdepth}{-\maxdimen} % remove section numbering
\usepackage{iftex}
\ifPDFTeX
  \usepackage[T1]{fontenc}
  \usepackage[utf8]{inputenc}
  \usepackage{textcomp} % provide euro and other symbols
\else % if luatex or xetex
  \usepackage{unicode-math} % this also loads fontspec
  \defaultfontfeatures{Scale=MatchLowercase}
  \defaultfontfeatures[\rmfamily]{Ligatures=TeX,Scale=1}
\fi
\usepackage{lmodern}
\ifPDFTeX\else
  % xetex/luatex font selection
  \setmainfont[]{Libertinus Serif}
  \setmonofont[Scale=0.85]{Latin Modern Mono}
  \setmathfont[]{Latin Modern Math}
\fi
% Use upquote if available, for straight quotes in verbatim environments
\IfFileExists{upquote.sty}{\usepackage{upquote}}{}
\IfFileExists{microtype.sty}{% use microtype if available
  \usepackage[]{microtype}
  \UseMicrotypeSet[protrusion]{basicmath} % disable protrusion for tt fonts
}{}
\makeatletter
\@ifundefined{KOMAClassName}{% if non-KOMA class
  \IfFileExists{parskip.sty}{%
    \usepackage{parskip}
  }{% else
    \setlength{\parindent}{0pt}
    \setlength{\parskip}{6pt plus 2pt minus 1pt}}
}{% if KOMA class
  \KOMAoptions{parskip=half}}
\makeatother
\usepackage{longtable,booktabs,array}
\newcounter{none} % for unnumbered tables
\usepackage{calc} % for calculating minipage widths
% Correct order of tables after \paragraph or \subparagraph
\usepackage{etoolbox}
\makeatletter
\patchcmd\longtable{\par}{\if@noskipsec\mbox{}\fi\par}{}{}
\makeatother
% Allow footnotes in longtable head/foot
\IfFileExists{footnotehyper.sty}{\usepackage{footnotehyper}}{\usepackage{footnote}}
\makesavenoteenv{longtable}
\setlength{\emergencystretch}{3em} % prevent overfull lines
\providecommand{\tightlist}{%
  \setlength{\itemsep}{0pt}\setlength{\parskip}{0pt}}

\providecommand{\E}{\mathbb{E}}
\providecommand{\N}{\mathbb{N}}
\providecommand{\Z}{\mathbb{Z}}
\providecommand{\R}{\mathbb{R}}
\providecommand{\eps}{\varepsilon}
\providecommand{\ceil}[1]{\left\lceil #1\right\rceil}
\providecommand{\floor}[1]{\left\lfloor #1\right\rfloor}
\AtBeginDocument{\let\setminus\smallsetminus}
\usepackage[most]{tcolorbox}
\newtcolorbox{warning}{colback=red!5,colframe=red!65!black,boxrule=0.8pt,arc=2pt,left=6pt,right=6pt,top=5pt,bottom=5pt}
\usepackage{etoolbox}
\AtBeginEnvironment{longtable}{\small}
\setlength{\emergencystretch}{3em}
\usepackage{fancyhdr}
\pagestyle{fancy}
\fancyhf{}
\fancyhead[L]{\small\bfseries UNREFEREED GPT-6 ASTRA OUTPUT}
\fancyhead[R]{\small\thepage}
\renewcommand{\headrulewidth}{0.4pt}
\usepackage{bookmark}
\IfFileExists{xurl.sty}{\usepackage{xurl}}{} % add URL line breaks if available
\urlstyle{same}
\hypersetup{
  pdftitle={Circular colouring the orthogonality graph},
  pdfauthor={Writeup: gpt-6-astra (single model pass)},
  colorlinks=true,
  linkcolor={blue!50!black},
  filecolor={Maroon},
  citecolor={Blue},
  urlcolor={blue!50!black},
  pdfcreator={LaTeX via pandoc}}

\title{Circular colouring the orthogonality graph}
\usepackage{etoolbox}
\makeatletter
\providecommand{\subtitle}[1]{% add subtitle to \maketitle
  \apptocmd{\@title}{\par {\large #1 \par}}{}{}
}
\makeatother
\subtitle{Unrefereed candidate proof by GPT-6 Astra}
\author{Writeup: \texttt{gpt-6-astra} (single model pass)}
\date{Catalog id \texttt{circular\_colouring\_the\_orthogonality\_graph}
--- generated 2026-09-16}

\begin{document}
\maketitle

\begin{warning}

\textbf{UNREFEREED MODEL OUTPUT.} This document was selected solely
because GPT-6 Astra labelled its own result
\texttt{would\_publish:\ true} and returned \texttt{proved} or
\texttt{disproved}. It has not passed the adversarial LLM referee used
for the earlier Sol campaign, has not been checked by a human
mathematician, and has not been checked for novelty. Treat every
mathematical and bibliographic claim below as unverified.

\end{warning}

\section{Summary and provenance}\label{summary-and-provenance}

{\def\LTcaptype{none} % do not increment counter
\begin{longtable}[]{@{}
  >{\raggedright\arraybackslash}p{(\linewidth - 2\tabcolsep) * \real{0.5000}}
  >{\raggedright\arraybackslash}p{(\linewidth - 2\tabcolsep) * \real{0.5000}}@{}}
\toprule\noalign{}
\endhead
\bottomrule\noalign{}
\endlastfoot
Catalog id & \texttt{circular\_colouring\_the\_orthogonality\_graph} \\
Catalog entry &
\href{https://graph-theory-ai.github.io/graph-conjectures/op/circular_colouring_the_orthogonality_graph/}{Circular
colouring the orthogonality graph} \\
Source corpus & opg \\
Campaign leg & selected second attempts (\texttt{attacks\_retry}) \\
Source paper / entry &
http://www.openproblemgarden.org/op/circular\_colouring\_the\_orthogonality\_graph \\
Model verdict & \textbf{proved} (confidence: high) \\
Model's one-line claim & A geometric odd-walk lemma excludes every
homomorphism from O to a finite locally bipartite graph, and
consequently proves χ\_c(O)=4. \\
Model & \texttt{gpt-6-astra}, reasoning effort \texttt{max},
\texttt{mode=pro}, flex service tier \\
Original artifact &
\texttt{attacks\_retry/circular\_colouring\_the\_orthogonality\_graph/output.md} \\
Independent review & \textbf{None} \\
Model's caveats & Literature priority has not been checked; the argument
below is self-contained. \\
\end{longtable}
}

\section{Problem statement}\label{problem-statement}

\emph{No extracted statement was stored in the artifact metadata.}

\section{Astra writeup}\label{astra-writeup}

\emph{The text below is the model output verbatim, apart from moving its
machine-readable verdict block into the summary above and shifting
Markdown heading levels for this document. Mechanical TeX defects and
equation tags were normalized where needed for compilation.}

\subsection{Circular chromatic number of the orthogonality
graph}\label{circular-chromatic-number-of-the-orthogonality-graph}

Let \(\mathcal O\) be the graph whose vertices are the one-dimensional
subspaces of \(\mathbb R^3\), with adjacency given by orthogonality.

We prove the following stronger statement.

\textbf{Theorem.} There is no homomorphism from \(\mathcal O\) to a
finite graph in which every open neighborhood induces a bipartite graph.
Consequently, \[
\boxed{\chi_c(\mathcal O)=4.}
\]

The key geometric fact is that the union of the neighborhoods of three
linearly independent directions always contains an odd closed walk.

\subsubsection{1. The geometric odd-walk
lemma}\label{the-geometric-odd-walk-lemma}

For a nonzero vector \(a\), write \[
P_a=a^\perp,\qquad J_a(x)=a\times x.
\] Thus \(J_a(x)\), when nonzero, represents a line perpendicular to
\([x]\) and belonging to \(P_a\).

\textbf{Lemma.} Let \(a,b,c\in\mathbb R^3\) be linearly independent. The
subgraph of \(\mathcal O\) induced by all lines contained in \[
P_a\cup P_b\cup P_c
\] is not bipartite.

\textbf{Proof.} First suppose that two of the normals, say \(a,b\), are
perpendicular. Then \[
[a],\quad [b],\quad [a\times b]
\] form a triangle, and all three belong to the specified union. This
covers that case.

We may therefore assume that \(a,b,c\) are pairwise nonorthogonal.
Normalize them to unit vectors and choose their signs and an orthonormal
coordinate system so that \[
a=(0,0,1),\qquad b=(s,0,t),\qquad c=(u,v,w),
\] where \[
s>0,\qquad 0<t<1,\qquad s^2+t^2=1.
\] Linear independence and pairwise nonorthogonality give \[
v\ne0,\qquad w\ne0,\qquad su+tw\ne0. \qquad\text{(1)}
\]

We will construct walks using successive cross products. All the cross
products involved are nonzero: whenever \(d\cdot e\ne0\), the
restriction \[
J_d:P_e\longrightarrow P_d
\] is injective, because the kernel of \(J_d\) is \(\mathbb R d\), which
does not meet \(P_e\) nontrivially. This applies to every transition
below.

Identify \(P_a\) with \(\mathbb R^2\) using
\((x,y,0)\leftrightarrow(x,y)\). Consider the two operators on \(P_a\)
\[
E=J_a^2J_b^2,\qquad F=J_aJ_cJ_b.
\] Direct calculation gives \[
E=
\begin{pmatrix}
t^2&0\\
0&1
\end{pmatrix},
\qquad
F=
\begin{pmatrix}
0&su+tw\\
-tw&sv
\end{pmatrix}. \qquad\text{(2)}
\]

These operators have a walk interpretation:

\begin{itemize}
\tightlist
\item
  \(F\) is a three-edge walk, with successive planes \[
  P_a\ \xrightarrow{J_b}\ P_b\
  \xrightarrow{J_c}\ P_c\
  \xrightarrow{J_a}\ P_a.
  \]
\item
  \(E\) is a four-edge walk, with successive planes \[
  P_a\ \xrightarrow{J_b}\ P_b\
  \xrightarrow{J_b}\ P_b\
  \xrightarrow{J_a}\ P_a\
  \xrightarrow{J_a}\ P_a.
  \]
\end{itemize}

For \(n\ge0\), \[
E^nF=
\begin{pmatrix}
0&t^{2n}(su+tw)\\
-tw&sv
\end{pmatrix}.
\] Its characteristic polynomial has discriminant \[
\Delta_n
=(sv)^2-4t^{2n}tw(su+tw). \qquad\text{(3)}
\] Because \(0<t<1\) and \(sv\ne0\), \[
\Delta_n\longrightarrow (sv)^2>0.
\] Choose \(n\ge0\) with \(\Delta_n>0\). Then \(E^nF\) has a real
eigenvector \(x_0\ne0\). Moreover, (1) shows that \(E^nF\) is
invertible, so its corresponding eigenvalue \(\lambda\) is nonzero: \[
E^nF(x_0)=\lambda x_0.
\]

Start at the line \([x_0]\), perform the three cross-product steps
defining \(F\), and then perform \(n\) copies of the four steps defining
\(E\). Every step joins perpendicular lines, every intermediate vector
is nonzero, and every line lies in \(P_a\cup P_b\cup P_c\). The final
line is \[
[E^nF(x_0)]=[\lambda x_0]=[x_0].
\] We have therefore obtained a closed walk of odd length \[
3+4n.
\] A bipartite graph cannot have an odd closed walk. This proves the
lemma. \(\square\)

A useful reformulation is the following. For
\(S\subseteq V(\mathcal O)\), let \[
N_{\mathcal O}(S)=\bigcup_{L\in S}N_{\mathcal O}(L).
\]

\textbf{Corollary.} If \(\mathcal O[N_{\mathcal O}(S)]\) is bipartite,
then the lines in \(S\) span a vector subspace of dimension at most
\(2\).

Indeed, otherwise \(S\) contains three linearly independent directions,
and the lemma supplies an odd closed walk inside \(N_{\mathcal O}(S)\).

\subsubsection{2. No finite locally bipartite
target}\label{no-finite-locally-bipartite-target}

Call a graph \emph{locally bipartite} if every open neighborhood induces
a bipartite graph.

Suppose, for a contradiction, that \[
h:\mathcal O\longrightarrow H
\] is a homomorphism to a finite locally bipartite graph \(H\).

For any \(\alpha\in V(H)\), put \[
S_\alpha=h^{-1}(\alpha).
\] Every vertex of \(N_{\mathcal O}(S_\alpha)\) is mapped into
\(N_H(\alpha)\). Consequently, \[
\mathcal O[N_{\mathcal O}(S_\alpha)]
\longrightarrow H[N_H(\alpha)].
\] The graph on the right is bipartite, so the graph on the left is
bipartite as well. By the corollary, every fiber \(S_\alpha\) spans a
subspace of dimension at most \(2\).

These finitely many fibers cannot cover all projective directions. Here
is an explicit pigeonhole argument. If \(m=|V(H)|\), consider the
\(2m+1\) lines \[
L_j=[(1,j,j^2)],\qquad 1\le j\le 2m+1.
\] Some three have the same image under \(h\). But any three distinct
such lines are linearly independent, since \[
\det
\begin{pmatrix}
1&i&i^2\\
1&j&j^2\\
1&k&k^2
\end{pmatrix}
=(j-i)(k-i)(k-j)\ne0.
\] This contradicts the dimension bound on that fiber.

Thus \(\mathcal O\) has no homomorphism to a finite locally bipartite
graph.

\subsubsection{\texorpdfstring{3. Circular colorings below \(4\) have
locally bipartite
targets}{3. Circular colorings below 4 have locally bipartite targets}}\label{circular-colorings-below-4-have-locally-bipartite-targets}

For positive integers \(p,q\) with \(p\ge2q\), let \(K_{p/q}\) be the
circular clique on \(\mathbb Z_p\), with adjacency defined by \[
a\sim b
\quad\Longleftrightarrow\quad
q\le (b-a\bmod p)\le p-q.
\]

\textbf{Lemma.} If \(p<4q\), then \(K_{p/q}\) is locally bipartite.

\textbf{Proof.} By translation symmetry, it is enough to consider the
neighborhood of \(0\): \[
N(0)=\{q,q+1,\ldots,p-q\}.
\] Partition it into \[
A=N(0)\cap\{q,\ldots,2q-1\},
\qquad
B=N(0)\cap\{2q,\ldots,3q-1\}.
\] These sets cover \(N(0)\), because \(p-q<3q\). Each lies in a block
of \(q\) consecutive colors, and such a block is independent in
\(K_{p/q}\). Thus \(A,B\) are a bipartition of the neighborhood.
\(\square\)

It follows from Section 2 that \[
\mathcal O\nrightarrow K_{p/q}
\qquad\text{whenever }p/q<4.
\] Using the finite-palette definition of circular chromatic number,
this already proves \[
\chi_c(\mathcal O)\ge4. \qquad\text{(4)}
\]

For completeness, the conclusion is unchanged if circular chromatic
number for infinite graphs is defined using real-circle colorings.
Suppose a circular \(r\)-coloring exists for some \(r<4\). Choose a
rational number \[
r<R<4
\] and scale the coloring to circumference \(R\). Every edge then has
circular distance at least \(R/r>1\). Write \(R=p/q\), choosing a
sufficiently large common multiple of its numerator and denominator that
\[
\frac{2}{q}<\frac Rr-1.
\] Rounding every color to the grid of spacing \(1/q\) changes each
color by less than \(1/q\), so every edge still has circular distance at
least \(1\). This gives a homomorphism to the finite graph \(K_{p/q}\),
a contradiction. The rounding estimate is uniform and therefore applies
to infinitely many vertices.

\subsubsection{4. A proper four-coloring, including boundary
directions}\label{a-proper-four-coloring-including-boundary-directions}

Here is an explicit four-coloring to establish the upper bound without
leaving boundary choices implicit.

For every line choose its unit representative \(u=(x,y,z)\) according to
the rule:

\begin{itemize}
\tightlist
\item
  choose \(z>0\), if possible;
\item
  if \(z=0\), choose \(y>0\), if possible;
\item
  if \(z=y=0\), choose \(x>0\).
\end{itemize}

Except for the line \([e_3]\), the planar projection \((x,y)\) is
nonzero. Color such a line by the half-open quadrant containing
\((x,y)\), using the argument intervals \[
[0,\pi/2),\quad [\pi/2,\pi),\quad
[\pi,3\pi/2),\quad [3\pi/2,2\pi).
\] Give \([e_3]\) the third color.

Two nonzero planar vectors in the same half-open quadrant have strictly
positive inner product. Since all selected representatives satisfy
\(z\ge0\), two non-pole lines of the same color have \[
u\cdot v
=(x,y)\cdot(x',y')+zz'>0.
\] Thus they are not perpendicular.

The selected representatives of equatorial lines have arguments in
\([0,\pi)\), so no equatorial line has the third color. Every other line
of the third color has \(z>0\) and hence is not perpendicular to
\(e_3\). Therefore the coloring is proper, and \[
\chi_c(\mathcal O)\le\chi(\mathcal O)\le4. \qquad\text{(5)}
\]

Combining (4) and (5) proves \[
\boxed{\chi_c(\mathcal O)=4}.
\]

\subsubsection{5. Finite witnesses and the previous
attempt}\label{finite-witnesses-and-the-previous-attempt}

The obstruction above can also be made into a fully specified finite
construction for each palette size.

Fix \(m\). Start with the \(2m+1\) lines \[
[(1,j,j^2)],\qquad 1\le j\le2m+1.
\] For each triple, apply the geometric lemma, choosing the least
\(n\ge0\) for which (3) is positive and an eigenline of the displayed
matrix \(E^nF\). Include all lines of its resulting odd closed walk. Let
\(F_m\) be the orthogonality graph induced on the resulting finite set
of lines.

Then \(F_m\) has no homomorphism to any locally bipartite graph with at
most \(m\) vertices: three of the initial lines would receive the same
color, and their associated odd walk would map into that color's
bipartite neighborhood. This construction uses only arithmetic, square
roots, and a terminating search for \(n\).

In particular, for each fixed circular clique \(K_{p/q}\) with
\(p/q<4\), there is a finite obstruction. The obstruction depends on the
palette size; the proof does \textbf{not} assert that one finite
subgraph has circular chromatic number \(4\). It is therefore compatible
with the previous attempt's finite-subgraph projection observation.

The proposed discrete-spectrum reduction is not needed here. The
decisive additional ingredient is the odd-walk lemma, which turns
bipartiteness of a color's neighborhood into a dimension restriction on
its entire fiber. No unproved conjecture, regularity assumption on the
coloring, or computational assertion is used.

\end{document}
